\subsubsection*{4.1 Error Correcting Codes}


\begin{beispiel}
Codewörter: \quad $\mathbf{c}_1 = 100110$, $\mathbf{c}_2 = 010101$, $\mathbf{c}_3 = 001011$, $\mathbf{c}_4 = 000000$
\end{beispiel}

Mit Wörtern wollen wir rechnen. Deshalb kommen die Buchstaben aus einem Körper $F = F_q$, \ $q$ Primzahlpotenz.

\begin{definition}[Code]
Der \textbf{Code} $C$ ist die Menge der Codewörter $C \subseteq F_q^{\,n}$, \ $n$ ist die Wortlänge.
\end{definition}

\begin{tabular}{ll}
obiges Beispiel: & $n = 6$, \ $q = 2$ \\
ISBN10: & $n = 10$, \ $q = 11$ \\
\end{tabular}

$q = 2$: binärer Code, \quad $q = 3$: ternärer Code.

\paragraph{Senden und Empfangen einer Nachricht $\mathbf{a} \in F_q^{\,k}$, \ $k < n$}
Codierungsfunktion \ $\varphi : F_q^{\,k} \to F_q^{\,n}$, \quad $\mathbf{a} = a_1 \dots a_k$, \quad $\varphi(\mathbf{a}) = \mathbf{c}$

% GRAFIK: Senkrechtes Ablaufdiagramm mit Pfeilen nach unten; beim Übertragungsschritt ein Pfeil von rechts "<- mögliche Fehler".
\[
\begin{array}{ll}
\mathbf{a} \text{ codieren: } \varphi(\mathbf{a}) = \mathbf{c} \in C \ \text{ Codewort} & \\
\quad\downarrow & \\
\mathbf{c} \text{ senden} & \\
\quad\downarrow & \qquad\leftarrow \quad \text{mögliche Fehler} \\
\mathbf{r} \in F_q^{\,n} \text{ empfangen} & \\
\quad\downarrow & \\
\mathbf{r} \text{ decodieren: } \mathbf{r} \to \mathbf{c} \in C \ \text{ gesendetes Codewort} & \\
\quad\downarrow & \\
\mathbf{a} = \varphi^{-1}(\mathbf{c}) \ \text{ Nachricht} &
\end{array}
\]

$\varphi$ injektiv (eineindeutig), aber nicht surjektiv.

Fehler \textbf{erkennen:} \ $\mathbf{r} \in C$ ?\\
Fehler \textbf{korrigieren:} \ $\mathbf{r} \to \mathbf{c}$. \ $\mathbf{r}$ auf das nächtgelegene Codewort $\mathbf{c}$ abbilden.

\paragraph{$\mathbf{c}$ ist unbekannt. Wie kann der Empfänger $\mathbf{r} \neq \mathbf{c}$ erkennen?}
\textbf{Lösung}: Es werden nicht alle Wörter aus $F_q^{\,n}$ als Codewörter zugelassen,
\[ C \neq F_q^{\,n} \]

\begin{beispiel}
$C = \{\mathbf{c}_1, \mathbf{c}_2, \mathbf{c}_3, \mathbf{c}_4\}$\\
Zwei Codewörter unterscheiden sich in mindestens 3 Stellen,\\
d.h.\ $e < 3$ Fehler werden erkannt.\\
$\mathbf{r} = 010001$ \ ist kein Codewort.\\
Unter der Annahme, dass 2 Fehler seltener sind als 1 Fehler, wird das Wort
\[ \mathbf{r} \to \mathbf{c}_2 = 010101 \]
rekonstruiert.
% KORRIGIERT: im Original stand "y" statt r (r = 010001 unterscheidet sich in 2 Stellen von c4 = 000000)
$\mathbf{r}$ unterscheidet sich in 2 Stellen von $\mathbf{c}_4$. 2 Fehler kann $C$ nicht korrigieren.
\end{beispiel}

% ERGÄNZT: Erklärung zum Beispiel (nicht im Original)
\begin{bemerkung}[Erklärung zum Beispiel]
$e$ = Anzahl der Fehler (gekippte Stellen), $\mathbf{r}$ = empfangenes Wort, $d(\mathbf{u}, \mathbf{v})$ = Anzahl verschiedener Stellen.
\begin{enumerate}
\item Minimalabstand $d_C = \min_{i \neq j} d(\mathbf{c}_i, \mathbf{c}_j) = 3$.
\item Erkennen: $1 \le e \le d_C - 1 = 2$ $\Rightarrow$ $\mathbf{r} \notin C$, denn ein anderes Codewort braucht $\ge 3$ Änderungen.
      Bei $e = 3$ kann $\mathbf{r}$ wieder ein Codewort sein.
\item Korrigieren = nächstes Codewort wählen. Abstände von $\mathbf{r} = 010001$:
\[
d(\mathbf{r}, \mathbf{c}_1) = 5, \quad d(\mathbf{r}, \mathbf{c}_2) = 1, \quad d(\mathbf{r}, \mathbf{c}_3) = 3, \quad d(\mathbf{r}, \mathbf{c}_4) = 2
\]
$\Rightarrow$ $\mathbf{r} \to \mathbf{c}_2$ (Unterschied nur an Stelle 4).
\item Wurde $\mathbf{c}_4$ gesendet ($e = 2$), ist diese Korrektur falsch: $\mathbf{c}_2$ liegt näher als $\mathbf{c}_4$.
\end{enumerate}
Merke: $C$ erkennt $d_C - 1 = 2$ Fehler und korrigiert $\lfloor (d_C - 1)/2 \rfloor = 1$ Fehler.
\end{bemerkung}

\begin{definition}[Gewicht]
Definieren das \textbf{Gewicht} eines Wortes $\mathbf{w} \in F_q^{\,n}$
\[ |\mathbf{w}| = |\{\, i \mid w_i \neq 0 \,\}| \]
\end{definition}

\begin{definition}[Hamming Abstand]
Der \textbf{Hamming Abstand} \quad $|\mathbf{v}-\mathbf{w}|$\\
gibt an, an wie vielen Stellen sich die Wörter $\mathbf{v}, \mathbf{w} \in F_q^{\,n}$ unterscheiden.
\end{definition}

Für den Hamming Abstand gelten die \textbf{Abstandsregeln}: \ $\mathbf{v}, \mathbf{w}, \mathbf{y} \in F_q^{\,n}$
\begin{center}
\begin{tabular}{ll}
$|\mathbf{v}-\mathbf{w}| = |\mathbf{w}-\mathbf{v}|$ & Symmetrie \\
$|\mathbf{v}-\mathbf{w}| \geq 0$ & positiv semidefinit \\
% KORRIGIERT: im Original stand "w = y" statt "v = w"
$|\mathbf{v}-\mathbf{w}| = 0 \ \Rightarrow \ \mathbf{v} = \mathbf{w}$ & positiv definit \\
$|\mathbf{v}-\mathbf{w}| \leq |\mathbf{v}-\mathbf{y}| + |\mathbf{y}-\mathbf{w}|$ & Dreiecksungleichung \\
\end{tabular}
\end{center}

\begin{definition}[Minimalgewicht]
\textbf{Minimalgewicht} des Codes $C$ \ (minimum distance)
\[ d_C := \min\{\, |\mathbf{c} - \mathbf{c}'| \mid \mathbf{c}, \mathbf{c}' \in C,\ \mathbf{c} \neq \mathbf{c}' \,\} \]
\end{definition}

\begin{beispiel}
Der obige Code $C$ hat das Minimalgewicht $d_C = 3$.
\end{beispiel}

\begin{satz}
% KORRIGIERT: im Original "Sei e = Anzahl der Fehler im Wort r ... C erkennt/korrigiert das Wort r"; die Äquivalenz gilt nur für "jedes Wort mit höchstens e Fehlern"
Sei $e \geq 1$ eine natürliche Zahl.\\
\begin{tabular}{lcl}
$C$ \textit{erkennt} jedes Wort $\mathbf{r}$ mit höchstens $e$ Fehlern als falsch & $\Leftrightarrow$ & $e + 1 \leq d_C$ \\
$C$ \textit{korrigiert} jedes Wort $\mathbf{r}$ mit höchstens $e$ Fehlern, $\mathbf{r} \to \mathbf{c}$ & $\Leftrightarrow$ & $2e + 1 \leq d_C$ \\
\end{tabular}
\end{satz}

\begin{definition}
Falls \ $d_C \geq e + 1$ \ heißt $C$ \textbf{$e$-Fehler-erkennend};\\
falls \ $d_C \geq 2e + 1$ \ heißt $C$ \textbf{$e$-Fehler-korrigierend}.
\end{definition}

\paragraph{Forderungen an einen Code $C$}
\begin{tabular}{ll}
- schnelles Codieren: & $\varphi(\mathbf{a}) = \mathbf{c}$ \\
- schnelles Decodieren: & $\mathbf{r} \to \mathbf{c}$, \ $\varphi^{-1}(\mathbf{c}) = \mathbf{a}$ \\
- viele Fehler korrigierend & $d_C$ groß, \\
- hohe Informationsdichte. & $|C| / |F_q^{\,n}|$ \\
\end{tabular}

\paragraph{Parameter des Codes $C$} \qquad $[n, |C|, d_C]_q$

% ------------------------------------------------------------
